% tuto_apic2d.tex -- Tutoriel : un solveur APIC 2D quasi-incompressible en Taichi
% Compiler avec pdfLaTeX (Overleaf : compilateur par défaut).
\documentclass[11pt,a4paper]{article}

\usepackage[utf8]{inputenc}
\usepackage[T1]{fontenc}
\usepackage{lmodern}
\usepackage[french]{babel}
\usepackage[margin=2.2cm]{geometry}
\usepackage{amsmath,amssymb}
\usepackage{xcolor}
\usepackage{booktabs}
\usepackage{tikz}
\usetikzlibrary{arrows.meta,positioning}
\usepackage{listings}
\usepackage{hyperref}
\hypersetup{colorlinks=true,linkcolor=blue!60!black,urlcolor=blue!60!black}

% ---------------------------------------------------------------- style des listings
\definecolor{codebg}{RGB}{247,247,250}
\definecolor{codekw}{RGB}{0,80,160}
\definecolor{codecomment}{RGB}{90,120,90}
\definecolor{codestring}{RGB}{160,60,20}
\lstset{
  language=Python,
  basicstyle=\ttfamily\small,
  keywordstyle=\color{codekw}\bfseries,
  commentstyle=\color{codecomment}\itshape,
  stringstyle=\color{codestring},
  backgroundcolor=\color{codebg},
  frame=single, framerule=0pt, rulecolor=\color{codebg},
  xleftmargin=6pt, xrightmargin=6pt,
  breaklines=true, breakatwhitespace=true,
  showstringspaces=false, tabsize=4,
  columns=fullflexible, keepspaces=true,
  literate=
    {é}{{\'e}}1 {è}{{\`e}}1 {ê}{{\^e}}1 {ë}{{\"e}}1
    {à}{{\`a}}1 {â}{{\^a}}1 {î}{{\^i}}1 {ï}{{\"i}}1
    {ô}{{\^o}}1 {ù}{{\`u}}1 {û}{{\^u}}1 {ç}{{\c{c}}}1
    {É}{{\'E}}1 {È}{{\`E}}1 {À}{{\`A}}1 {œ}{{\oe}}1
}
% Code en ligne : \detokenize neutralise les caractères actifs de babel-french (: ; ! ?)
% et affiche les underscores tels quels, sans passer par \lstinline (fragile dans une macro).
\newcommand{\code}[1]{\texttt{\detokenize{#1}}}

\title{Un solveur APIC 2D d'eau quasi-incompressible en Taichi\\
  \large Tutoriel pas à pas, de la page blanche au programme complet}
\author{}
\date{}

\begin{document}
\maketitle

\begin{abstract}
Ce tutoriel construit, en six étapes et environ 150 lignes de Python/Taichi, un solveur de fluide
2D hybride particules/grille de type \textbf{APIC} (\emph{Affine Particle-In-Cell}).
L'eau est traitée comme un fluide \textbf{quasi-incompressible} : la pression vient d'une simple
loi d'état, sans résolution d'équation de Poisson. Le programme gère les parois du domaine,
des obstacles, et affiche au choix le champ de masse volumique ou de vitesse dans la fenêtre
GGUI de Taichi. Le code complet, testé avec Taichi 1.7.4, est donné en annexe.
\end{abstract}

\tableofcontents
\bigskip

% ================================================================
\section{Prérequis}

\begin{itemize}
  \item Python 3.10 à 3.13 et Taichi :
\begin{lstlisting}[language=bash]
pip install taichi
\end{lstlisting}
  \item Un GPU est un plus (CUDA ou Vulkan) mais pas indispensable : remplacer
        \code{ti.gpu} par \code{ti.cpu} dans \code{ti.init} suffit.
  \item Lancer le programme final avec \code{python apic2d.py}. Touches :
        \textbf{Espace} bascule masse volumique / vitesse, \textbf{R} réinitialise, \textbf{Échap} quitte.
\end{itemize}

Le tutoriel a été testé avec Taichi 1.7.4 (Python 3.13, backend CUDA et backend CPU).

% ================================================================
\section{L'idée générale}

\subsection{Une méthode hybride}

Une méthode \emph{particules/grille} stocke le fluide sur des \textbf{particules} (position, vitesse,
état de compression) et utilise une \textbf{grille} temporaire à chaque pas de temps pour calculer
les forces et appliquer les conditions limites. Les particules transportent la matière
(l'advection est exacte et sans diffusion numérique), la grille fait le travail eulérien
(pression, parois, obstacles). Chaque pas de temps est un cycle en trois temps :

\begin{center}
\begin{tikzpicture}[
  box/.style={draw, rounded corners, minimum width=3.6cm, minimum height=1.1cm, align=center, fill=blue!5},
  arr/.style={-{Latex[length=3mm]}, thick}]
  \node[box] (p) {Particules\\ $x_p,\,v_p,\,C_p,\,J_p$};
  \node[box, right=3.4cm of p] (g) {Grille\\ $m_i,\,v_i$};
  \draw[arr] (p.north east) to[bend left=25] node[above]{\textbf{1. P2G} : masse, quantité de mouvement, pression} (g.north west);
  \draw[arr] (g.south west) to[bend left=25] node[below]{\textbf{3. G2P} : vitesse, $C_p$, $J_p$, position} (p.south east);
  \draw[arr] (g.east) -- ++(1.6,0) node[right, align=left]{\textbf{2. Grille} :\\ gravité,\\ parois, obstacles};
\end{tikzpicture}
\end{center}

\begin{enumerate}
  \item \textbf{P2G} (\emph{particle to grid}) : chaque particule dépose sa masse et sa quantité
        de mouvement sur les n\oe uds voisins, ainsi que la force de pression qu'elle exerce.
  \item \textbf{Mise à jour de la grille} : on obtient la vitesse des n\oe uds, on ajoute la gravité,
        on impose les conditions limites (parois, obstacles).
  \item \textbf{G2P} (\emph{grid to particle}) : chaque particule relit la vitesse de la grille,
        met à jour sa matrice affine $C_p$, son taux de compression $J_p$ et avance.
\end{enumerate}

\subsection{PIC, FLIP, APIC}

La différence entre les variantes tient à ce qu'une particule transporte entre deux pas :
\begin{itemize}
  \item \textbf{PIC} : la particule prend la vitesse interpolée de la grille. Simple mais très
        dissipatif : chaque aller-retour lisse le champ de vitesse.
  \item \textbf{FLIP} : la particule ajoute l'\emph{incrément} de vitesse de la grille. Peu dissipatif
        mais bruité, il faut le mélanger avec du PIC.
  \item \textbf{APIC} : la particule transporte sa vitesse \emph{et} une matrice $C_p$ ($2\times 2$)
        décrivant la variation locale de vitesse autour d'elle. On dépose sur la grille une vitesse
        \emph{affine} $v_p + C_p (x_i - x_p)$. Résultat : conservation du moment cinétique, peu
        de dissipation, pas de bruit, aucun paramètre de mélange. C'est ce que nous implémentons.
\end{itemize}

\subsection{Quasi-incompressible : pourquoi et comment}

Un solveur strictement incompressible impose $\nabla\cdot v = 0$ en résolvant une équation de
Poisson pour la pression à chaque pas (\emph{projection}). C'est la partie la plus lourde d'un
solveur de fluide. Ici on l'évite : chaque particule suit son \textbf{taux de dilatation}
$J_p$ ($J_p = 1$ au repos, $J_p < 1$ comprimée) et la pression est donnée par une loi d'état
\[
  p = E\,(1 - J),
\]
où $E$ est une raideur. Le fluide est un peu compressible, avec une vitesse du son
$c_s = \sqrt{E/\rho}$. Tant que les vitesses de l'écoulement restent petites devant $c_s$
(nombre de Mach faible), les variations de masse volumique restent de l'ordre de $\mathrm{Ma}^2$,
soit quelques pour cent : c'est \og quasi-incompressible \fg. Le prix à payer est un pas de temps
petit, limité par $c_s$ (voir \S\ref{sec:stab}).

% ================================================================
\section{Le minimum de mathématiques}

\subsection{Interpolation par B-spline quadratique}
\label{sec:bspline}

La grille a $n$ n\oe uds par côté sur le carré unité, espacés de $\Delta x = 1/n$.
Une particule en $x_p$ interagit avec les $3\times 3$ n\oe uds les plus proches via le noyau
B-spline quadratique (par dimension) :
\[
  N(u) =
  \begin{cases}
    \tfrac34 - u^2 & |u| < \tfrac12,\\[2pt]
    \tfrac12\left(\tfrac32 - |u|\right)^2 & \tfrac12 \le |u| < \tfrac32,\\[2pt]
    0 & \text{sinon.}
  \end{cases}
\]
Dans le code on calcule le n\oe ud \code{base} en bas à gauche du stencil et la position
relative \code{fx} $\in [0.5, 1.5)^2$ :
\[
  \texttt{base} = \left\lfloor \frac{x_p}{\Delta x} - \tfrac12 \right\rfloor,\qquad
  \texttt{fx} = \frac{x_p}{\Delta x} - \texttt{base},
\]
et les trois poids par dimension valent
\[
  w_0 = \tfrac12\,(1.5 - \texttt{fx})^2,\qquad
  w_1 = \tfrac34 - (\texttt{fx} - 1)^2,\qquad
  w_2 = \tfrac12\,(\texttt{fx} - 0.5)^2,
\]
avec $w_0 + w_1 + w_2 = 1$. Le poids 2D du n\oe ud $(i,j)$ du stencil est $w_{ip} = w_i^x\, w_j^y$.

\begin{center}
\begin{tikzpicture}[scale=1.1]
  \draw[step=1, gray!50, thin] (-0.5,-0.5) grid (3.5,3.5);
  \foreach \i in {0,...,3} \foreach \j in {0,...,3} \fill[gray!60] (\i,\j) circle (2pt);
  \foreach \i in {0,1,2} \foreach \j in {0,1,2} \fill[blue!70] (\i,\j) circle (3pt);
  \fill[red] (1.3,0.8) circle (3.5pt) node[right=3pt] {$x_p$};
  \node[below left, blue!70!black] at (0,0) {\texttt{base}};
  \draw[{Latex}-{Latex}, thick] (0,-0.9) -- (1.3,-0.9) node[midway, below] {$\texttt{fx}_x\,\Delta x$};
  \node[align=left, anchor=west] at (4.2,1.5) {n\oe uds du stencil $3\times3$\\ (bleus), pondérés par $w_{ip}$};
\end{tikzpicture}
\end{center}

\subsection{Le transfert APIC}

\paragraph{P2G.} Masse et quantité de mouvement des n\oe uds :
\[
  m_i = \sum_p w_{ip}\, m_p,\qquad
  (m v)_i = \sum_p w_{ip}\, m_p \big( v_p + C_p\,(x_i - x_p) \big).
\]
Le terme $C_p (x_i - x_p)$ est la signature d'APIC : la particule dépose un champ de vitesse
affine, pas une vitesse constante.

\paragraph{G2P.} Vitesse et matrice affine de la particule :
\[
  v_p = \sum_i w_{ip}\, v_i,\qquad
  C_p = B_p D_p^{-1},\quad
  B_p = \sum_i w_{ip}\, v_i\,(x_i - x_p)^{\mathsf T},\quad
  D_p = \sum_i w_{ip}\,(x_i - x_p)(x_i - x_p)^{\mathsf T}.
\]
Pour la B-spline quadratique on a exactement $D_p = \frac{\Delta x^2}{4}\, I$, d'où la formule
utilisée dans le code :
\[
  \boxed{\;C_p = \frac{4}{\Delta x^2}\sum_i w_{ip}\, v_i\,(x_i - x_p)^{\mathsf T}\;}
\]
$C_p$ est une approximation du gradient de vitesse $\nabla v$ au voisinage de la particule.

\subsection{Pression quasi-incompressible}

\paragraph{Taux de dilatation.} $J_p$ suit $\dot J = J\,\nabla\cdot v$, et comme
$C_p \approx \nabla v$, on a $\nabla\cdot v \approx \operatorname{tr} C_p$ :
\[
  J_p \leftarrow J_p\,\big(1 + \Delta t\,\operatorname{tr} C_p\big).
\]

\paragraph{Force nodale.} Avec la contrainte de Cauchy $\sigma = -p\,I = E(J-1)\,I$, la force
exercée par la particule $p$ sur le n\oe ud $i$ est $f_i = -V_p\,\sigma_p\,\nabla w_{ip}$.
Pour la B-spline quadratique le gradient du poids se simplifie en
$\nabla w_{ip} \approx \frac{4}{\Delta x^2}\,w_{ip}\,(x_i - x_p)$ (approximation MLS-MPM), donc
\[
  \Delta t\, f_i = w_{ip}\;\underbrace{\Big(-\Delta t\,\frac{4 E V_p}{\Delta x^2}\,(J_p - 1)\Big)}_{\texttt{stress}}\;(x_i - x_p).
\]
C'est exactement le terme \code{stress * dpos} déposé dans P2G. Vérification du signe :
une particule comprimée ($J_p < 1$) donne \code{stress} $> 0$ et pousse les n\oe uds
\emph{loin} d'elle, ce qui la détend.

\subsection{Stabilité et choix des paramètres}
\label{sec:stab}

Le schéma est explicite : il faut $\Delta t < \Delta x / (c_s + |v|_{\max})$ avec
$c_s = \sqrt{E/\rho}$. Avec les valeurs du code :

\begin{center}
\begin{tabular}{@{}lll@{}}
\toprule
Paramètre & Valeur & Rôle \\
\midrule
$n$ & 128 & n\oe uds par côté, $\Delta x = 1/128$ \\
$\rho$ & 1 & masse volumique de repos \\
$E$ & 400 & raideur, $c_s = 20$ \\
$\Delta t$ & $2\cdot10^{-4}$ & $< \Delta x / c_s \approx 3.9\cdot10^{-4}$ \\
$V_p$ & $(\Delta x/2)^2$ & 4 particules par cellule \\
$g$ & 9.8 & gravité \\
\code{bound} & 3 & épaisseur des parois (en n\oe uds) \\
\bottomrule
\end{tabular}
\end{center}

Une vitesse d'éclaboussure de $3$ donne $\mathrm{Ma} \approx 0.15$, donc des variations de masse
volumique de l'ordre de $2\,\%$. Augmenter $E$ rend l'eau plus incompressible mais impose
de réduire $\Delta t$.

% ================================================================
\section{Étape 1 : squelette Taichi, paramètres et champs}

On initialise Taichi, on définit les constantes physiques, puis les \emph{champs} (tableaux
vivant sur le GPU). Les particules ont quatre attributs, la grille en a deux plus un masque
d'obstacles, et une image sert à l'affichage.

\begin{lstlisting}
import taichi as ti

ti.init(arch=ti.gpu)  # remplacer par ti.cpu si pas de GPU

# ---------------------------------------------------------------- Paramètres
n_grid = 128                 # noeuds par côté (domaine = carré unité)
dx = 1.0 / n_grid
inv_dx = float(n_grid)
dt = 2e-4                    # pas de temps (doit vérifier dt < dx / c_s)
substeps = 10                # sous-pas par image affichée
bound = 3                    # épaisseur des parois du domaine (en noeuds)

rho = 1.0                    # masse volumique de repos
E = 400.0                    # raideur : pression p = E (1 - J), c_s = sqrt(E / rho)
gravity = 9.8
p_vol = (dx * 0.5) ** 2      # 4 particules par cellule -> volume d'une particule
p_mass = p_vol * rho

# Bloc d'eau initial [x0, x1] x [y0, y1]
water_x0, water_x1 = 0.05, 0.45
water_y0, water_y1 = 0.30, 0.90
n_particles = int((water_x1 - water_x0) * (water_y1 - water_y0) * n_grid * n_grid * 4)

# Obstacles : un disque et un rectangle
disk_center = ti.Vector([0.55, 0.30])
disk_radius = 0.10
rect_x0, rect_x1, rect_y0, rect_y1 = 0.72, 0.95, 0.00, 0.14

# ---------------------------------------------------------------- Champs
x = ti.Vector.field(2, ti.f32, n_particles)          # position
v = ti.Vector.field(2, ti.f32, n_particles)          # vitesse
C = ti.Matrix.field(2, 2, ti.f32, n_particles)       # matrice affine APIC (gradient de vitesse)
J = ti.field(ti.f32, n_particles)                    # taux de dilatation (1 = repos)

grid_v = ti.Vector.field(2, ti.f32, (n_grid, n_grid))  # quantité de mouvement puis vitesse
grid_m = ti.field(ti.f32, (n_grid, n_grid))            # masse
solid = ti.field(ti.i32, (n_grid, n_grid))             # 1 si le noeud est dans un obstacle
img = ti.Vector.field(3, ti.f32, (n_grid, n_grid))     # image affichée
\end{lstlisting}

\paragraph{Remarques.}
\begin{itemize}
  \item \code{grid_v} contient la \emph{quantité de mouvement} après P2G, puis la \emph{vitesse}
        après division par la masse : on réutilise le même champ.
  \item \code{n_particles} est calculé pour avoir 4 particules par cellule dans le bloc d'eau,
        cohérent avec \code{p_vol}.
  \item Le masque \code{solid} est défini sur les n\oe uds : un n\oe ud dans un obstacle est
        marqué 1. C'est tout ce qu'il faut pour gérer des obstacles de forme quelconque.
\end{itemize}

% ================================================================
\section{Étape 2 : initialisation (eau et obstacles)}

Un kernel place les particules aléatoirement dans le bloc d'eau, au repos ($v = 0$,
$C = 0$, $J = 1$), et remplit le masque d'obstacles par des tests géométriques sur la
position de chaque n\oe ud. Le même kernel sert de \emph{reset} (touche R).

\begin{lstlisting}
@ti.kernel
def init():
    for p in x:
        x[p] = [water_x0 + ti.random() * (water_x1 - water_x0),
                water_y0 + ti.random() * (water_y1 - water_y0)]
        v[p] = [0.0, 0.0]
        C[p] = ti.Matrix.zero(ti.f32, 2, 2)
        J[p] = 1.0
    for i, j in solid:
        pos = ti.Vector([i, j]) * dx
        in_disk = (pos - disk_center).norm() < disk_radius
        in_rect = rect_x0 <= pos.x and pos.x <= rect_x1 and rect_y0 <= pos.y and pos.y <= rect_y1
        solid[i, j] = 1 if (in_disk or in_rect) else 0
\end{lstlisting}

Pour ajouter un obstacle, il suffit d'ajouter un test sur \code{pos} : n'importe quelle forme
décrite analytiquement (ou lue depuis une image) convient.

% ================================================================
\section{Étape 3 : P2G, des particules vers la grille}

C'est le c\oe ur du solveur. Pour chaque particule : stencil, poids, coefficient de pression,
puis dépôt sur les 9 n\oe uds. Les opérateurs \code{+=} sur des champs sont automatiquement
atomiques dans une boucle parallèle Taichi.

\begin{lstlisting}
@ti.kernel
def p2g():
    for i, j in grid_m:
        grid_v[i, j] = [0.0, 0.0]
        grid_m[i, j] = 0.0
    for p in x:
        base = int(x[p] * inv_dx - 0.5)          # noeud en bas à gauche du stencil 3x3
        fx = x[p] * inv_dx - base                # position relative dans le stencil
        w = [0.5 * (1.5 - fx) ** 2,              # poids B-spline quadratique
             0.75 - (fx - 1.0) ** 2,
             0.5 * (fx - 0.5) ** 2]
        # Pression quasi-incompressible : p = E (1 - J). Le coefficient "stress"
        # contient dt, le volume, et le 4/dx^2 du gradient des poids.
        stress = -dt * 4.0 * E * p_vol * (J[p] - 1.0) * inv_dx * inv_dx
        for i, j in ti.static(ti.ndrange(3, 3)):
            offset = ti.Vector([i, j])
            dpos = (offset - fx) * dx            # vecteur particule -> noeud
            weight = w[i].x * w[j].y
            grid_v[base + offset] += weight * (p_mass * (v[p] + C[p] @ dpos)  # APIC
                                               + stress * dpos)               # pression
            grid_m[base + offset] += weight * p_mass
\end{lstlisting}

\paragraph{Ligne par ligne.}
\begin{itemize}
  \item \code{base}, \code{fx}, \code{w} : voir \S\ref{sec:bspline}. \code{fx} et chaque
        \code{w[k]} sont des vecteurs 2D ; \code{w[i].x * w[j].y} est le poids 2D du n\oe ud
        $(i,j)$ du stencil.
  \item \code{stress} : le coefficient $-\Delta t\,4 E V_p (J_p-1)/\Delta x^2$ de la force de pression.
  \item \code{ti.static(ti.ndrange(3, 3))} : boucle $3\times3$ déroulée à la compilation.
  \item \code{dpos} $= x_i - x_p$ en unités physiques.
  \item Le dépôt combine le terme APIC $m_p (v_p + C_p\,\texttt{dpos})$ et le terme de pression
        $\texttt{stress}\cdot\texttt{dpos}$, tous deux pondérés par $w_{ip}$.
\end{itemize}

% ================================================================
\section{Étape 4 : mise à jour de la grille, parois et obstacles}

Sur chaque n\oe ud portant de la masse : conversion en vitesse, gravité, puis conditions limites.

\begin{lstlisting}
@ti.kernel
def grid_step():
    for i, j in grid_m:
        if grid_m[i, j] > 0:
            grid_v[i, j] /= grid_m[i, j]                      # quantité de mouvement -> vitesse
            grid_v[i, j].y -= dt * gravity
            # Parois du domaine (glissement libre) : on annule la composante normale entrante
            if i < bound and grid_v[i, j].x < 0:
                grid_v[i, j].x = 0.0
            if i > n_grid - bound and grid_v[i, j].x > 0:
                grid_v[i, j].x = 0.0
            if j < bound and grid_v[i, j].y < 0:
                grid_v[i, j].y = 0.0
            if j > n_grid - bound and grid_v[i, j].y > 0:
                grid_v[i, j].y = 0.0
            # Obstacles (adhérence) : vitesse nulle
            if solid[i, j] == 1:
                grid_v[i, j] = [0.0, 0.0]
\end{lstlisting}

\paragraph{Parois du domaine.} Sur une bande de \code{bound} n\oe uds le long de chaque bord,
on annule la composante de vitesse \emph{normale} qui pointe vers la paroi, et on laisse la
composante tangentielle : c'est une condition de \textbf{glissement libre} (\emph{free-slip})
qui autorise l'eau à se décoller de la paroi. Les particules ne peuvent pas traverser :
une fois dans la bande, tous les n\oe uds de leur stencil interdisent la vitesse entrante.

\paragraph{Obstacles.} Sur les n\oe uds marqués \code{solid}, la vitesse est mise à zéro :
condition d'\textbf{adhérence} (\emph{no-slip}), la plus simple possible. L'eau colle un peu à
l'obstacle et les particules peuvent s'y enfoncer d'une cellule environ (cosmétique : la zone
est dessinée en gris). Une variante glissement libre demanderait la normale de l'obstacle en
chaque n\oe ud.

% ================================================================
\section{Étape 5 : G2P, de la grille vers les particules}

Chaque particule relit la vitesse des 9 n\oe uds, reconstruit $C_p$ avec la formule encadrée,
met à jour $J_p$ et avance.

\begin{lstlisting}
@ti.kernel
def g2p():
    for p in x:
        base = int(x[p] * inv_dx - 0.5)
        fx = x[p] * inv_dx - base
        w = [0.5 * (1.5 - fx) ** 2,
             0.75 - (fx - 1.0) ** 2,
             0.5 * (fx - 0.5) ** 2]
        new_v = ti.Vector.zero(ti.f32, 2)
        new_C = ti.Matrix.zero(ti.f32, 2, 2)
        for i, j in ti.static(ti.ndrange(3, 3)):
            offset = ti.Vector([i, j])
            dpos = (offset - fx) * dx
            g_v = grid_v[base + offset]
            weight = w[i].x * w[j].y
            new_v += weight * g_v
            new_C += 4.0 * inv_dx * inv_dx * weight * g_v.outer_product(dpos)  # C = sum w v (x_i - x_p)^T / (dx^2/4)
        v[p] = new_v
        C[p] = new_C
        J[p] *= 1.0 + dt * new_C.trace()         # dJ/dt = J div(v)
        x[p] += dt * v[p]
        x[p] = ti.math.clamp(x[p], bound * dx, 1.0 - bound * dx)  # sécurité : jamais hors grille
\end{lstlisting}

\paragraph{Remarques.}
\begin{itemize}
  \item \code{g_v.outer_product(dpos)} calcule $v_i\,(x_i-x_p)^{\mathsf T}$ ; le facteur
        \code{4.0 * inv_dx * inv_dx} est le $D_p^{-1} = 4/\Delta x^2$.
  \item \code{new_C.trace()} $\approx \nabla\cdot v$ sert à mettre à jour $J_p$.
  \item Le \code{clamp} final est une sécurité : une particule hors de la grille produirait
        un accès mémoire invalide dans P2G (silencieux sur GPU). Avec les parois de l'étape 4,
        il n'est normalement jamais actif.
\end{itemize}

À ce stade le solveur est complet : un pas de temps est \code{p2g(); grid_step(); g2p()}.

% ================================================================
\section{Étape 6 : affichage et boucle principale}

\subsection{Le kernel de rendu}

On remplit une image de $n\times n$ pixels directement depuis la grille : gris pour les
obstacles, sinon une couleur proportionnelle à la masse volumique (mode 0) ou à la norme de la
vitesse (mode 1). La masse volumique d'un n\oe ud est $m_i/\Delta x^2$, ramenée à $\rho$
(valeur 1 pour de l'eau au repos). Un léger gain d'affichage (1/0.8) sature le volume d'eau
pour que l'\oe il ne soit pas distrait par le bruit d'échantillonnage des particules.

\begin{lstlisting}
@ti.kernel
def render(mode: ti.i32):
    for i, j in img:
        if solid[i, j] == 1:
            img[i, j] = [0.35, 0.35, 0.35]
        else:
            val = 0.0
            if mode == 0:
                val = grid_m[i, j] / (dx * dx * rho) / 0.8   # masse volumique relative (1 = eau au repos), gain 1/0.8
            else:
                val = grid_v[i, j].norm() / 3.0        # vitesse (3 m/s -> couleur maximale)
            val = ti.math.clamp(val, 0.0, 1.0)
            if mode == 0:
                img[i, j] = ti.Vector([0.02, 0.02, 0.08]) * (1 - val) + ti.Vector([0.35, 0.65, 1.0]) * val
            else:
                img[i, j] = ti.Vector([0.02, 0.02, 0.08]) * (1 - val) + ti.Vector([1.0, 0.75, 0.2]) * val
\end{lstlisting}

\subsection{La fenêtre GGUI et la boucle}

\code{ti.ui.Window} ouvre une fenêtre ; \code{canvas.set_image} affiche un champ image et le met
automatiquement à l'échelle de la fenêtre ; \code{window.get_gui()} fournit un petit panneau de
texte. Les événements clavier se lisent avec \code{window.get_event(ti.ui.PRESS)} qui remplit
\code{window.event.key}. À chaque image on effectue \code{substeps} pas de temps.

\begin{lstlisting}
def main():
    window = ti.ui.Window("APIC 2D - eau quasi-incompressible", res=(512, 512))
    canvas = window.get_canvas()
    gui = window.get_gui()
    mode = 0
    init()
    while window.running:
        while window.get_event(ti.ui.PRESS):
            if window.event.key == ti.ui.ESCAPE:
                window.running = False
            elif window.event.key == ti.ui.SPACE:
                mode = 1 - mode
            elif window.event.key == 'r':
                init()
        for _ in range(substeps):
            p2g()
            grid_step()
            g2p()
        render(mode)
        canvas.set_image(img)
        gui.begin("APIC 2D", 0.02, 0.02, 0.50, 0.16)
        gui.text("Affichage : " + ("masse volumique" if mode == 0 else "vitesse"))
        gui.text("ESPACE : basculer l'affichage")
        gui.text("R : reset      ESC : quitter")
        gui.end()
        window.show()


if __name__ == "__main__":
    main()
\end{lstlisting}

\subsection{Ce que l'on doit voir}

Le bloc d'eau tombe, s'écrase sur le sol, contourne le disque, submerge le rectangle puis
s'étale. En mode masse volumique, les zones d'impact apparaissent plus claires (compression) ;
en mode vitesse, les jets et les éclaboussures ressortent en orange. Avec ces paramètres,
la masse totale sur la grille est conservée exactement, $J_p$ reste entre $0.9$ et $1.1$ pour
plus de 99\,\% des particules, et la simulation reste stable indéfiniment.

% ================================================================
\section{Pour aller plus loin}

Trois extensions naturelles, dans l'ordre de difficulté : (1) une condition de glissement
libre sur les obstacles, en stockant une normale par n\oe ud solide et en n'annulant que la
composante normale entrante ; (2) une vraie \textbf{projection de pression} (résolution de
$\nabla^2 p = \rho\,\nabla\cdot v / \Delta t$ sur la grille par Jacobi ou gradient conjugué),
qui remplace $E$ et $J$ et lève la contrainte sur $\Delta t$ liée à $c_s$ ; (3) le passage en 3D,
quasi immédiat : vecteurs à trois composantes, stencil $3\times3\times3$ et $D_p$ inchangé.

% ================================================================
\appendix
\section{Code complet : \texttt{apic2d.py}}

\begin{lstlisting}
# apic2d.py -- Solveur APIC 2D pour de l'eau quasi-incompressible (Taichi)
#
# Cycle par pas de temps :
#   1. p2g()       : particules -> grille (masse, quantité de mouvement, force de pression)
#   2. grid_step() : grille (vitesse, gravité, parois, obstacles)
#   3. g2p()       : grille -> particules (vitesse, matrice affine C, J, position)
#
# Touches : ESPACE = masse volumique / vitesse, R = reset, ESC = quitter

import taichi as ti

ti.init(arch=ti.gpu)  # remplacer par ti.cpu si pas de GPU

# ---------------------------------------------------------------- Paramètres
n_grid = 128                 # noeuds par côté (domaine = carré unité)
dx = 1.0 / n_grid
inv_dx = float(n_grid)
dt = 2e-4                    # pas de temps (doit vérifier dt < dx / c_s)
substeps = 10                # sous-pas par image affichée
bound = 3                    # épaisseur des parois du domaine (en noeuds)

rho = 1.0                    # masse volumique de repos
E = 400.0                    # raideur : pression p = E (1 - J), c_s = sqrt(E / rho)
gravity = 9.8
p_vol = (dx * 0.5) ** 2      # 4 particules par cellule -> volume d'une particule
p_mass = p_vol * rho

# Bloc d'eau initial [x0, x1] x [y0, y1]
water_x0, water_x1 = 0.05, 0.45
water_y0, water_y1 = 0.30, 0.90
n_particles = int((water_x1 - water_x0) * (water_y1 - water_y0) * n_grid * n_grid * 4)

# Obstacles : un disque et un rectangle
disk_center = ti.Vector([0.55, 0.30])
disk_radius = 0.10
rect_x0, rect_x1, rect_y0, rect_y1 = 0.72, 0.95, 0.00, 0.14

# ---------------------------------------------------------------- Champs
x = ti.Vector.field(2, ti.f32, n_particles)          # position
v = ti.Vector.field(2, ti.f32, n_particles)          # vitesse
C = ti.Matrix.field(2, 2, ti.f32, n_particles)       # matrice affine APIC (gradient de vitesse)
J = ti.field(ti.f32, n_particles)                    # taux de dilatation (1 = repos)

grid_v = ti.Vector.field(2, ti.f32, (n_grid, n_grid))  # quantité de mouvement puis vitesse
grid_m = ti.field(ti.f32, (n_grid, n_grid))            # masse
solid = ti.field(ti.i32, (n_grid, n_grid))             # 1 si le noeud est dans un obstacle
img = ti.Vector.field(3, ti.f32, (n_grid, n_grid))     # image affichée


# ---------------------------------------------------------------- Initialisation
@ti.kernel
def init():
    for p in x:
        x[p] = [water_x0 + ti.random() * (water_x1 - water_x0),
                water_y0 + ti.random() * (water_y1 - water_y0)]
        v[p] = [0.0, 0.0]
        C[p] = ti.Matrix.zero(ti.f32, 2, 2)
        J[p] = 1.0
    for i, j in solid:
        pos = ti.Vector([i, j]) * dx
        in_disk = (pos - disk_center).norm() < disk_radius
        in_rect = rect_x0 <= pos.x and pos.x <= rect_x1 and rect_y0 <= pos.y and pos.y <= rect_y1
        solid[i, j] = 1 if (in_disk or in_rect) else 0


# ---------------------------------------------------------------- 1. Particules -> grille
@ti.kernel
def p2g():
    for i, j in grid_m:
        grid_v[i, j] = [0.0, 0.0]
        grid_m[i, j] = 0.0
    for p in x:
        base = int(x[p] * inv_dx - 0.5)          # noeud en bas à gauche du stencil 3x3
        fx = x[p] * inv_dx - base                # position relative dans le stencil
        w = [0.5 * (1.5 - fx) ** 2,              # poids B-spline quadratique
             0.75 - (fx - 1.0) ** 2,
             0.5 * (fx - 0.5) ** 2]
        # Pression quasi-incompressible : p = E (1 - J). Le coefficient "stress"
        # contient dt, le volume, et le 4/dx^2 du gradient des poids.
        stress = -dt * 4.0 * E * p_vol * (J[p] - 1.0) * inv_dx * inv_dx
        for i, j in ti.static(ti.ndrange(3, 3)):
            offset = ti.Vector([i, j])
            dpos = (offset - fx) * dx            # vecteur particule -> noeud
            weight = w[i].x * w[j].y
            grid_v[base + offset] += weight * (p_mass * (v[p] + C[p] @ dpos)  # APIC
                                               + stress * dpos)               # pression
            grid_m[base + offset] += weight * p_mass


# ---------------------------------------------------------------- 2. Mise à jour de la grille
@ti.kernel
def grid_step():
    for i, j in grid_m:
        if grid_m[i, j] > 0:
            grid_v[i, j] /= grid_m[i, j]                      # quantité de mouvement -> vitesse
            grid_v[i, j].y -= dt * gravity
            # Parois du domaine (glissement libre) : on annule la composante normale entrante
            if i < bound and grid_v[i, j].x < 0:
                grid_v[i, j].x = 0.0
            if i > n_grid - bound and grid_v[i, j].x > 0:
                grid_v[i, j].x = 0.0
            if j < bound and grid_v[i, j].y < 0:
                grid_v[i, j].y = 0.0
            if j > n_grid - bound and grid_v[i, j].y > 0:
                grid_v[i, j].y = 0.0
            # Obstacles (adhérence) : vitesse nulle
            if solid[i, j] == 1:
                grid_v[i, j] = [0.0, 0.0]


# ---------------------------------------------------------------- 3. Grille -> particules
@ti.kernel
def g2p():
    for p in x:
        base = int(x[p] * inv_dx - 0.5)
        fx = x[p] * inv_dx - base
        w = [0.5 * (1.5 - fx) ** 2,
             0.75 - (fx - 1.0) ** 2,
             0.5 * (fx - 0.5) ** 2]
        new_v = ti.Vector.zero(ti.f32, 2)
        new_C = ti.Matrix.zero(ti.f32, 2, 2)
        for i, j in ti.static(ti.ndrange(3, 3)):
            offset = ti.Vector([i, j])
            dpos = (offset - fx) * dx
            g_v = grid_v[base + offset]
            weight = w[i].x * w[j].y
            new_v += weight * g_v
            new_C += 4.0 * inv_dx * inv_dx * weight * g_v.outer_product(dpos)  # C = sum w v (x_i - x_p)^T / (dx^2/4)
        v[p] = new_v
        C[p] = new_C
        J[p] *= 1.0 + dt * new_C.trace()         # dJ/dt = J div(v)
        x[p] += dt * v[p]
        x[p] = ti.math.clamp(x[p], bound * dx, 1.0 - bound * dx)  # sécurité : jamais hors grille


# ---------------------------------------------------------------- Affichage
@ti.kernel
def render(mode: ti.i32):
    for i, j in img:
        if solid[i, j] == 1:
            img[i, j] = [0.35, 0.35, 0.35]
        else:
            val = 0.0
            if mode == 0:
                val = grid_m[i, j] / (dx * dx * rho) / 0.8   # masse volumique relative (1 = eau au repos), gain 1/0.8
            else:
                val = grid_v[i, j].norm() / 3.0        # vitesse (3 m/s -> couleur maximale)
            val = ti.math.clamp(val, 0.0, 1.0)
            if mode == 0:
                img[i, j] = ti.Vector([0.02, 0.02, 0.08]) * (1 - val) + ti.Vector([0.35, 0.65, 1.0]) * val
            else:
                img[i, j] = ti.Vector([0.02, 0.02, 0.08]) * (1 - val) + ti.Vector([1.0, 0.75, 0.2]) * val


# ---------------------------------------------------------------- Boucle principale
def main():
    window = ti.ui.Window("APIC 2D - eau quasi-incompressible", res=(512, 512))
    canvas = window.get_canvas()
    gui = window.get_gui()
    mode = 0
    init()
    while window.running:
        while window.get_event(ti.ui.PRESS):
            if window.event.key == ti.ui.ESCAPE:
                window.running = False
            elif window.event.key == ti.ui.SPACE:
                mode = 1 - mode
            elif window.event.key == 'r':
                init()
        for _ in range(substeps):
            p2g()
            grid_step()
            g2p()
        render(mode)
        canvas.set_image(img)
        gui.begin("APIC 2D", 0.02, 0.02, 0.50, 0.16)
        gui.text("Affichage : " + ("masse volumique" if mode == 0 else "vitesse"))
        gui.text("ESPACE : basculer l'affichage")
        gui.text("R : reset      ESC : quitter")
        gui.end()
        window.show()


if __name__ == "__main__":
    main()
\end{lstlisting}

\end{document}
